\section{Exact counting and the shifted product}\label{cms:sec:counting}
\subsection{The complete finite counting polynomial}
For a permutation $\pi\in S_n$, write $c(\pi)$ for the sum of its cycle minima. Define
\[
G_n(z)=\sum_{\pi\in S_n}z^{c(\pi)},\qquad G_0(z)=1.
\]
\begin{proposition}\label{cms:prop:finite}
For every integer $n\geq1$,
\begin{equation}\label{cms:eq:finite}
G_n(z)=\prod_{k=1}^{n}(k-1+z^k).
\end{equation}
The same polynomial enumerates the sum of record positions in permutations of $[n]$. Consequently $a_n=[z^n]G_n(z)$ for $n\geq0$.
\end{proposition}
\begin{proof}
Insert the largest label $n$ into a permutation of $[n-1]$ written in cycles. There are $n-1$ possible insertions immediately after one of the existing labels. Each leaves every cycle minimum unchanged. The remaining possibility creates the singleton cycle $(n)$, increasing the sum of minima by $n$. Every permutation of $[n]$ is obtained exactly once in this way. Thus
\[
G_n(z)=(n-1+z^n)G_{n-1}(z),
\]
and iteration proves \eqref{cms:eq:finite}.

For the second interpretation, let $I_k$ indicate that position $k$ is a left-to-right maximum of a uniform random permutation. The relative rank of the entry in position $k$ among the first $k$ entries is uniform on $\{1,\ldots,k\}$, and these relative ranks are independent. One way to see independence is that the relative-rank code is a bijection between permutations and the $1\cdot2\cdots n=n!$ possible codes. Therefore $I_k$ are independent with $\Pr(I_k=1)=1/k$, and
\[
\mathbb E z^{\sum_{k=1}^{n}kI_k}
=\prod_{k=1}^{n}\left(1-\frac1k+\frac{z^k}{k}\right)
=\frac{G_n(z)}{n!}.
\]
This proves equality in distribution of the two statistics and hence the equality of their counting polynomials. It is not a claim that the two sums agree for each individual permutation.
\end{proof}

Factoring out the first factor $z$ and the constants in the other factors gives
\begin{equation}\label{cms:eq:finitefactor}
G_n(z)=z(n-1)!\prod_{j=1}^{n-1}\left(1+\frac{z^{j+1}}j\right)
\qquad(n\geq1).
\end{equation}
Thus
\begin{equation}\label{cms:eq:shift}
b_n=[z^{n-1}]\prod_{j=1}^{n-1}\left(1+\frac{z^{j+1}}j\right)
=[z^{n-1}]F(z).
\end{equation}
The factor with $j=n-1$ already has degree $n$ and is harmless for this coefficient; the same is true of every subsequent factor. The equality is exact for $n=1$ as well: the constant coefficient is $1$. At $n=2$ the coefficient of $z$ vanishes, explaining the exceptional zero $a_2=0$.

\subsection{Analytic existence and a global growth bound}
For each $0\leq r<1$,
\[
\sum_{j\geq1}\frac{r^{j+1}}j=-r\log(1-r)<\infty.
\]
The product \eqref{cms:eq:Fintro} therefore converges locally uniformly on $\D=\{z:|z|<1\}$. It defines an analytic function with nonnegative Taylor coefficients. None of the factors vanishes in the disk, since its zeros have modulus $j^{1/(j+1)}\geq1$. Hence $F$ is nonzero in $\D$. Moreover,
\begin{equation}\label{cms:eq:globalbound}
|F(z)|\leq F(|z|)
\leq\exp\!\left(\sum_{j\geq1}\frac{|z|^{j+1}}j\right)
=(1-|z|)^{-|z|}\leq\frac1{1-|z|}.
\end{equation}
This is the finite global order $-1$ used in the classical hybrid theorem. It also controls the function independently of any local expansion.

\subsection{The leading equivalent was already implied by a local limit theorem}
Put $T_n=\sum_{k=1}^{n}kI_k$. Proposition~\ref{cms:prop:finite} gives
\[
\Pr(T_n=n)=\frac{a_n}{n!},\qquad n\Pr(T_n=n)=b_n.
\]
Theorem~2.1 of \cite{GSW} states that, for integer $\kappa_n$ with $\kappa_n/n\to x>0$,
\[
n\Pr(T_n=\kappa_n)\longrightarrow\ee^{-\gamma}\rho(x),
\]
where the Dickman function satisfies $\rho(x)=1$ on $[0,1]$. Taking $\kappa_n=n$ gives $b_n\to A$. This report also computes $A$ directly from the product, because that residue enters the global subtraction. Neither the alternative derivation nor its use here makes the leading equivalent a new result.
